\newpage
\section{Background: The \Apeiron{} Framework and Multi-Axial Atomicity}

\subsection{The \Apeiron{} Framework}

The \Apeiron{} Framework~\cite{beniers2026apeiron} is a seventeen-layer architecture for autonomous knowledge genesis, designed to operate without inheriting human categories, temporalities, or causal assumptions. The framework is organized into four clusters:
\begin{itemize}
    \item \textbf{Cluster I (Substrate \& Emergence, Layers~1--4):} Establishes irreducible observables, relational hypergraphs, functional clusters, and endogenous dynamics.
    \item \textbf{Cluster II (Adaptive Intelligence, Layers~5--7):} Enables optimization, meta-learning, and emergent self-awareness.
    \item \textbf{Cluster III (Ontological Fluidity, Layers~8--13):} Supports the creation of new categories, reconciliation of incommensurable ontologies, and ontogenesis of novelty.
    \item \textbf{Cluster IV (World-Making \& Governance, Layers~14--17):} Facilitates autopoietic worldbuilding, responsible morphogenesis, collective cognition, and absolute integration.
\end{itemize}

This paper focuses exclusively on \textbf{Layer~1: Foundational Observables (The Epistemic Singularity)}, which defines the minimal, irreducible units of representation that ground the entire architecture.

\subsection{Multi-Axial Atomicity}

The \Apeiron{} Framework posits that an observable is \textit{atomic} (i.e., cannot be further decomposed without loss of essential identity or generative capacity) if it is simultaneously recognized as such by four independent mathematical frameworks:

\begin{enumerate}
    \item \textbf{Boolean Atomicity (Proposition~A.1 in~\cite{beniers2026apeiron}):} An element $a$ in a Boolean algebra is an atom if $a \neq 0$ and there exists no $b$ with $0 < b < a$. For integers under divisibility, the only atom is $1$. For sets under inclusion, atoms are the singleton sets.
    \item \textbf{Measure-Theoretic Atomicity (Proposition~A.2):} A measurable set $A$ is an atom for a measure $\mu$ if $\mu(A) > 0$ and every measurable subset $B \subseteq A$ satisfies either $\mu(B) = 0$ or $\mu(B) = \mu(A)$. In practice, this is approximated by checking subsets for the counting measure.
    \item \textbf{Categorical Atomicity (Proposition~A.3):} An object in a category is a \textit{zero object} (and hence atomic) if it is both initial (unique morphism from it to every other object) and terminal (unique morphism from every other object to it). For pointed sets, the singleton set is a zero object.
    \item \textbf{Information-Theoretic Atomicity (Proposition~A.4):} An observable is atomic if its Kolmogorov complexity $K(x)$ is approximately equal to its length $|x|$; i.e., it is algorithmically incompressible. In practice, $K(x)$ is bounded above by the size of the zlib-compressed representation.
\end{enumerate}

An observable that satisfies all four criteria is considered \textit{multi-axially atomic} and serves as a secure foundational unit for higher layers.